\documentclass[12pt,reqno]{amsart}
\usepackage{amsmath,amssymb,amsthm}
\usepackage{mathptmx}
\usepackage[scaled=0.90]{helvet}
\usepackage{courier}
\usepackage[T1]{fontenc}
\usepackage{enumitem}
\usepackage[hidelinks]{hyperref}
\pagestyle{plain}

\setlength{\oddsidemargin}{0cm}
\setlength{\evensidemargin}{0cm}
\setlength{\topmargin}{0in}
\setlength{\headheight}{0pt}
\setlength{\headsep}{0pt}
\setlength{\footskip}{.3in}
\setlength{\textheight}{9.2in}
\setlength{\textwidth}{6.5in}
\setlength{\parindent}{0pt}
\setlist[enumerate,1]{leftmargin=*,itemsep=1.1em}
\setlist[enumerate,2]{label=(\alph*),leftmargin=1.7em,itemsep=.45em,topsep=.4em}

\newcommand{\C}{\mathbb{C}}
\newcommand{\R}{\mathbb{R}}

\title{Problem Set 4: Solutions \\ MATH 411: Honors Complex Variables, Fall 2026}
\date{last compiled: \today}

% Coverage: Laurent series and Laurent expansions (Lectures 21--22),
% ahead of Midterm 2 on Fri, Nov. 6. Problems on isolated singularities,
% residues, the residue theorem, evaluating integrals, the argument
% principle, Rouche's theorem, and the open mapping theorem (Lectures 23--25)
% were moved to Problem Set 5 on Nov. 1, 2026.
% Suggested completion date: Wednesday, November 4, 2026.
% Sources (paths relative to the course workspace):
% 1: moved from Problem Set 3, problem 7
%    (originally Old Homework/Homework/Homework5_Fall2021.tex, 3--4).
% 2--3: Old Homework/Homework/Homework6_Fall2021.tex, 1 and 4 (hints
%    added).

\usepackage[most]{tcolorbox}
\definecolor{solblue}{RGB}{232,243,252}
\definecolor{solframe}{RGB}{120,170,215}
\newenvironment{solution}
  {\par\smallskip\begin{tcolorbox}[breakable,enhanced,colback=solblue,colframe=solframe,
   boxrule=0.5pt,arc=2pt,left=6pt,right=6pt,top=4pt,bottom=4pt]\setlength{\abovedisplayskip}{7pt}\setlength{\belowdisplayskip}{7pt}\noindent\textbf{Solution.}}
  {\hfill$\square$\end{tcolorbox}\par\medskip}
\usepackage{needspace}
\widowpenalty=10000
\clubpenalty=10000
\raggedbottom

% Solution 1 follows Old Homework/Homework/hwk3.tex, problems 7--8,
% and Homework5_Fall2021solutions.pdf, problems 3--4; the intermediate
% exterior-series typo is corrected and convergence is made explicit.
% Solution 2's alternative follows Homework6_Fall2021solutions.tex, 1,
% correcting the contour radius, product inequality, and equality case.
% The main proof of 2 follows the current power-sum hint without Rouche.
% Solution 3 follows the handwritten Homework 6 Solutions.pdf (2021), pp. 3--4.

\begin{document}
\maketitle

Unless stated otherwise, closed curves are oriented counterclockwise.

\begin{enumerate}

\item \textbf{Laurent expansions.}
\begin{enumerate}
\item Find a series in integer powers of $z$ representing
\[
 f(z)=\frac{1}{(z-1)(z-2)}
\]
in each of the regions $|z|<1$, $1<|z|<2$, and $|z|>2$.
Justify convergence in each region.
\item Find the first four nonzero terms of the Laurent expansion of
\[
 \frac{e^z}{z(1+z^2)}
\]
valid for $0<|z|<1$.
\end{enumerate}


\begin{solution}
\noindent (a) Note that
\[
\frac{1}{(z-1)(z-2)}=\frac{1}{z-2}-\frac{1}{z-1}.
\]
For $|z|<1$, expand both fractions in nonnegative powers of $z$:
\[
\boxed{f(z)=-\frac12\sum_{m=0}^\infty\left(\frac z2\right)^m
+\sum_{m=0}^\infty z^m
=\sum_{m=0}^\infty\left(1-\frac{1}{2^{m+1}}\right)z^m.}
\]
For $1<|z|<2$, write $-1/(z-1)=-z^{-1}/(1-1/z)$; the previous
expansion of $1/(z-2)$ still holds. Thus
\[
\boxed{f(z)=-\sum_{m=0}^\infty\frac{1}{z^{m+1}}
-\sum_{m=0}^\infty\frac{z^m}{2^{m+1}},\qquad 1<|z|<2.}
\]
For $|z|>2$, the same expansion of $-1/(z-1)$ holds, but now
\[
\frac{1}{z-2}=\frac1z\frac{1}{1-2/z}
=\sum_{m=0}^\infty\frac{2^m}{z^{m+1}}.
\]
Subtracting the two series gives
\[
\boxed{f(z)=\sum_{m=0}^\infty\frac{2^m-1}{z^{m+1}},\qquad |z|>2.}
\]
In the three regions the respective geometric ratios are $z,z/2$;
$1/z,z/2$; and $1/z,2/z$. Their moduli are strictly less than $1$.
All series therefore converge absolutely in the stated regions and
uniformly on compact subsets, justifying the additions and subtractions.

\medskip
\noindent (b) Multiply the expansions:
\begin{align*}
\frac{e^z}{z(1+z^2)}
&=\frac1z\left(1+z+\frac{z^2}{2!}+\frac{z^3}{3!}+\cdots\right)
          \left(1-z^2+z^4-\cdots\right)\\
&=\boxed{\frac1z+1-\frac z2-\frac56z^2+O(z^3).}
\end{align*}
These are the first four nonzero terms. The exponential series converges
absolutely everywhere, and the geometric series does so for $|z|<1$.
Their product is therefore valid for $0<|z|<1$.
\end{solution}

\needspace{12\baselineskip}
\item \textbf{Roots of a trinomial.}
Let $a\in\C$ and let $n\geq2$ be an integer. Prove that the equation
\[
 1+z+az^n=0
\]
has at least one root in the closed disc $|z|\leq2$.
\emph{Hint:} For $a\neq0$, the reciprocals $w_j=1/z_j$ of the roots are
the roots of $w^n+w^{n-1}+a$. Compute the power sums
$\sum_j w_j^{k}$ for $1\leq k\leq n-1$.


\begin{solution}
If $a=0$, take $z=-1$. Otherwise, let $z_1,\ldots,z_n$ be the nonzero
roots, counted with multiplicity, and put $w_j=1/z_j$.
Factoring $P(t)=1+t+at^n$ and comparing constant terms gives
\[
P(t)=\prod_{j=1}^n(1-w_jt).
\]
Taking the logarithmic derivative and expanding for sufficiently small
$|t|$, we find
\[
-\frac{tP'(t)}{P(t)}
=\sum_{j=1}^n\frac{w_jt}{1-w_jt}
=\sum_{k=1}^\infty\left(\sum_{j=1}^n w_j^k\right)t^k.
\]
On the other hand,
\[
-\frac{tP'(t)}{P(t)}
=-\frac{t+nat^n}{1+t+at^n}
=-\frac{t}{1+t}+O(t^n)
=\sum_{k=1}^{n-1}(-1)^k t^k+O(t^n).
\]
Comparing coefficients through degree $n-1$ gives
\[
\boxed{\sum_{j=1}^n w_j^k=(-1)^k\qquad(1\leq k\leq n-1).}
\]
If all $|z_j|>2$, then all $|w_j|<1/2$. Taking $k=n-1$ gives
\[
1=\left|\sum_{j=1}^n w_j^{n-1}\right|
\leq\sum_{j=1}^n|w_j|^{n-1}
<\frac{n}{2^{n-1}}\leq1,
\]
a contradiction. Here $n\leq2^{n-1}$ for $n\geq2$, as follows by
induction from $2=2^1$ and $n+1\leq2n$.
Thus at least one root satisfies $|z_j|\leq2$.
\end{solution}

\needspace{20\baselineskip}
\begin{solution}
\textit{Alternative from the previous homework solutions, using
Rouch\'e's theorem (a later topic).}
Write $F(z)=1+z+az^n$ and $G(z)=z+az^n$.
If $|a|<2^{-n}$, then on $|z|=2$,
\[
|G(z)|=2|1+az^{n-1}|
\geq2(1-|a|2^{n-1})>1=|F(z)-G(z)|.
\]
Rouch\'e's theorem implies that $F$ and $G$ have the same number of
zeros in $|z|<2$. Since $G(0)=0$, $F$ has at least one such zero.

If $|a|\geq2^{-n}$, factor
$F(z)=a(z-z_1)\cdots(z-z_n)$. Comparing constant terms gives
\[
|z_1|\cdots|z_n|=\frac1{|a|}\leq2^n.
\]
The roots cannot all have modulus greater than $2$, which proves the
claim, including the case $|a|=2^{-n}$.
\end{solution}

\needspace{12\baselineskip}
\item \textbf{No holomorphic extension of $1/z$.}
Show that there is no function $f$ that is holomorphic on the open unit
disc, extends continuously to the closed unit disc, and satisfies
$f(z)=1/z$ for $|z|=1$.


\begin{solution}
Suppose such a function $f$ existed and put $g(z)=zf(z)$. Then $g$
is holomorphic on the open disc and continuous on the closed disc.
At the centre $g(0)=0$, whereas on the unit circle $g(z)=1$.
By the mean value property, for every $0<r<1$,
\[
 0=g(0)=\frac{1}{2\pi}\int_0^{2\pi}g(re^{it})\,dt.
\]
We claim that the same formula holds at $r=1$. Indeed, $g$ is
uniformly continuous on the closed disc, and $|re^{it}-e^{it}|=1-r$
for every $t$. Thus
\[
\left|\frac{1}{2\pi}\int_0^{2\pi}
   \bigl(g(re^{it})-g(e^{it})\bigr)\,dt\right|
\leq \sup_t |g(re^{it})-g(e^{it})| \longrightarrow 0
\]
as $r\uparrow1$. But the boundary average is $1$, contradicting $g(0)=0$.
\end{solution}

\end{enumerate}
\end{document}
